\lecture{Macchine Operatrici}{07~-~10~-~2026}
%
\section{Pompe di circolazione}
% La pompa centrifuga è essenzialmente costituita da una
% girante che ruota eccentrica in una camera chiusa
% (chiocciola), aspirando liquido dalla parte centrale e
% inviandolo in posizione periferica dove è poi raccolto
% in una tubazione. La forma del profilo della girante,
% curvato all’indietro, serve per mantenere il fluido
% aderente alla superficie e nel contempo offrire una
% sezione di passaggio crescente, che consente di ridurre
% la velocità in uscita (l’energia cinetica si converte in
% pressione)
Vengono utilizzate per vincere le cadute di pressione e le variazioni positive di quota in un impianto idraulico.

Le pompe utilizzate in un impianto termotecnico sono tipicamente centrifughe.
%
\paragraph{Centrifughe} Sfruttano la forza centrifuga per concentrare il liquido nella parte esterna di un dispositivo a geometria cilindrica, incrementandone così la pressione e richiamando altro fluido nella zona centrale, dove la pressione è inferiore.
%
\subsection{Prevalenza}
Un'importante grandezza per il dimensionamento delle pompe è la \emphgreen{prevalenza}.

Richiamiamo l'equazione di conservazione dell'energia meccanica, scritta in termini di pressione:
\[
    p_{1} + \rho \frac{v^{2}_{1}}{2} + \rho \cdot g \cdot z_{1} + \Delta p_{t} = p_{2} + \rho \frac{v_{2}^{2}}{2} + \rho \cdot g \cdot z_{2} + \Delta p_{a}
\]
Tutti i termini che compongono questa equazione sono di fatto delle pressioni. La pompa di circolazione è chiamata a vincere tutte le differenze di pressione che si verificano tra l'ingresso e l'uscita:
\[
    \Delta p_{t} = (p_{2}-p_{1}) + (\rho \cdot \frac{v_{2}^{2}}{2}-\rho \cdot \frac{v_{1}^{2}}{2} ) + (\rho \cdot g \cdot z_{2} - \rho \cdot g \cdot z_{1})+ \Delta p_{a}
\]
in cui i termini a secondo membro rappresentano:
\begin{itemize}
    \item variazione di pressione statica: $ (p_{2}-p_{1}) $
    \item variazione di pressione cinetica: $ (\rho \cdot \frac{v_{2}^{2}}{2}-\rho \cdot \frac{v_{1}^{2}}{2} ) $
    \item variazione di quota potenziale (gravitazionale, anche detta geodetica): $ (\rho \cdot g \cdot z_{2} - \rho \cdot g \cdot z_{1}) $
    \item perdite di pressione per l'attrito: $ \Delta p_{a} $
\end{itemize}
Le variazione di energia cinetica risulta normalmente trascurabile, quindi possiamo semplificare l'equazione in:
\[
    \Delta p_{t} = (p_{2}-p_{1}) + \rho \cdot g \cdot (z_{2}-z_{1}) + \Delta p_{a}
\]
La valutazione in termini di quota del contributo legato alla pompa di circolazione si definisce \emphgreen{prevalenza}, viene indicata come $ H $ e si ottiene dalla precedente relazione dividendo tutti i termini per la densità e l'accelerazione di gravità:
\begin{formula}[Prevalenza]
    \[
        H = \frac{(p_{2}-p_{1})}{\rho \cdot g} + (z_{2}-z_{1}) + \frac{\Delta p_{a}}{\rho \cdot g}
    \]
\end{formula}
\begin{itemize}
    \item Il termine $ \frac{(p_{2}-p_{1})}{\rho \cdot g} $ si annulla nel caso in cui la pressione nella sezione iniziale e finale siano uguali (es. circuito idraulico in cui viene prelevata acqua da un recipiente $ A $ aperto e viene inviata ad un recipiente $ B $ aperto: entrambe le superfici sono soggette alla pressione atmosferica, la differenza di pressione statica è nulla);
    \item il termine $ (z_{2}-z_{1}) $ è nullo nel caso di un impianto idraulico \textbf{orizzontale}. Il termine è anche complessivamente nullo nel caso di un \emph{circuito chiuso}, in quanto la sezione iniziale e finale, che coincidono, si trovano alla stessa quota;
\end{itemize}

%
\subsection{Potenza assorbita dalla pompa}
